%% SYNTHETIC TEST UNIT generated by writing/tools/make_outline_book.py from FINAL_OUTLINE_v2 -- not book text
\chapter{横动量依赖（TMD）因子化}\label{ch:tmd}
本单元为测试用的合成文本（由 FINAL\_OUTLINE\_v2 自动生成），陈述内容从略。

\section{第 1 节（测试）}\label{sec:tmd-kinematics}

\begin{definition}[小 $q_T$ 区域]\label{def:tmd-kin}
（测试用合成陈述：内容从略。）
\uses{def:dy-kin}
\end{definition}

\begin{proposition}[小 $q_T$ 的领头区域 \tagA]\label{prop:tmd-regions}
（测试用合成陈述：内容从略。）
\uses{thm:ir-leading-regions, prop:ir-power-counting, def:tmd-kin}
\end{proposition}
\begin{proofsketch}测试。\end{proofsketch}
\paragraph{地位与缺口}测试：本陈述的地位由\cref{asm:appN-envelope}、\cref{asm:appN-separation}决定。

\begin{definition}[未减除的 TMD 关联函数]\label{def:tmd-unsub}
（测试用合成陈述：内容从略。）
\uses{def:prelim-wilson, def:pdf-quark}
\end{definition}

\begin{definition}[软因子]\label{def:tmd-soft}
（测试用合成陈述：内容从略。）
\uses{def:prelim-wilson}
\end{definition}

\section{第 2 节（测试）}\label{sec:tmd-regions}

\begin{lemma}[单圈软因子 \tagR]\label{lem:tmd-soft-oneloop}
（$R_{\mathrm{PT}}$：在重整化微扰论中对红外正规化的部分子态严格。）
（测试用合成陈述：内容从略。）
\uses{def:tmd-soft, prop:prelim-wilson-rules, def:sudakov-cusp}
\end{lemma}

\begin{lemma}[平方根构造是恒等式 \tagR]\label{lem:tmd-sqrt-identity}
（$R_{\mathrm{def}}$：由算符定义直接推出。）
（测试用合成陈述：内容从略。）
\uses{def:tmd-unsub, def:tmd-soft}
\end{lemma}

\begin{definition}[Collins 的 TMD 定义]\label{def:tmd-collins}
（测试用合成陈述：内容从略。）
\uses{def:tmd-unsub, def:tmd-soft, lem:tmd-sqrt-identity}
\end{definition}

\begin{remark}[其他方案]\label{rem:tmd-schemes}
（测试用合成陈述：内容从略。）
% uses: none
\end{remark}

\section{第 3 节（测试）}\label{sec:tmd-definitions}

\begin{remark}[方案常数（$S_\varepsilon$ 约定）]\label{rem:tmd-scheme-constants}
（测试用合成陈述：内容从略。）
% uses: none
\end{remark}

\begin{remark}[极限次序]\label{rem:tmd-order-of-limits}
（测试用合成陈述：内容从略。）
% uses: none
\end{remark}

\begin{proposition}[零箱减除即软因子 \tagR]\label{prop:tmd-zero-bin-soft}
（$R_{\mathrm{cond}}$：在所列编号假设之下证明完整；结论的地位不高于其中最弱的假设。）
前提：\cref{asm:appN-ward-product}、\cref{asm:appN-tmd}。
（测试用合成陈述：内容从略。）
\uses{def:tmd-collins, def:method-subtraction, lem:tmd-soft-oneloop, asm:appN-ward-product, asm:appN-tmd, lem:appN-absorption}
\end{proposition}

\begin{remark}[快度窗口]\label{rem:tmd-rapidity-window}
（测试用合成陈述：内容从略。）
% uses: none
\end{remark}

\section{第 4 节（测试）}\label{sec:tmd-subtraction}

\begin{lemma}[TMD DY、SIDIS、$e^+e^-$ 中的 Glauber 相消 \tagA]\label{lem:tmd-glauber}
（测试用合成陈述：内容从略。）
\uses{lem:dy-fsi, prop:appO-tmd, asm:appO-mixing, asm:appO-deformation, asm:appO-rewritten-cs}
\end{lemma}
\begin{proofsketch}测试。\end{proofsketch}
\paragraph{地位与缺口}测试：本陈述的地位由\cref{asm:appO-mixing}、\cref{asm:appO-deformation}、\cref{asm:appO-rewritten-cs}、\cref{rem:appO-literature-comparison}决定。

\begin{proposition}[小 $q_T$ DY 对模板假设的验证 \tagA]\label{prop:tmd-h1h5}
（测试用合成陈述：内容从略。）
\uses{thm:method-template, prop:tmd-regions, lem:tmd-glauber, prop:tmd-zero-bin-soft, lem:tmd-sqrt-identity}
\end{proposition}
\begin{proofsketch}测试。\end{proofsketch}
\paragraph{地位与缺口}测试：本陈述的地位由\cref{asm:appN-tmd}、\cref{asm:appN-ward-product}、\cref{asm:appO}决定。

\begin{theorem}[Drell–Yan 的 TMD 因子化 \tagA]\label{thm:tmd-factorization}
（测试用合成陈述：内容从略。）
非微扰地位：\cref{asm:intro-bridge}。
\uses{prop:tmd-h1h5, prop:tmd-regions, lem:tmd-soft-oneloop, lem:tmd-glauber, def:tmd-collins, thm:method-template, thm:dy-factorization, asm:intro-bridge, asm:appN-tmd, asm:appN-collins, asm:appO-rewritten-cs, asm:appO-mixing}
\end{theorem}
\begin{proofsketch}测试。\end{proofsketch}
\paragraph{地位与缺口}测试：本陈述的地位由\cref{asm:appN-tmd}、\cref{asm:appN-collins}、\cref{asm:appO-rewritten-cs}、\cref{asm:appO-mixing}、\cref{asm:intro-bridge}决定。

\begin{proposition}[Collins–Soper 方程与重整化群方程 \tagA]\label{prop:tmd-evolution}
（测试用合成陈述：内容从略。）
\uses{def:tmd-collins, lem:tmd-sqrt-identity, asm:appN-tmd}
\end{proposition}
\begin{proofsketch}测试。\end{proofsketch}
\paragraph{地位与缺口}测试：本陈述的地位由\cref{asm:appN-tmd}、\cref{prop:lattice-cs-kernel}决定。

\section{第 5 节（测试）}\label{sec:tmd-theorem}

\begin{proposition}[单圈演化系数与硬因子的 RG 相容性 \tagR]\label{prop:tmd-evolution-oneloop}
（$R_{\mathrm{PT}}$：在重整化微扰论中对红外正规化的部分子态严格。）
（测试用合成陈述：内容从略。）
\uses{lem:tmd-soft-oneloop, def:tmd-collins, prop:sudakov-oneloop}
\end{proposition}

\begin{proposition}[小 $b$ 展开 \tagA]\label{prop:tmd-small-b}
（测试用合成陈述：内容从略。）
非微扰地位：\cref{asm:intro-composite}、\cref{asm:intro-bridge}。
\uses{def:tmd-collins, thm:pdf-dglap, asm:intro-composite, asm:intro-bridge}
\end{proposition}
\begin{proofsketch}测试。\end{proofsketch}
\paragraph{地位与缺口}测试：见文献。

\begin{proposition}[CSS 解与非微扰输入 \tagA]\label{prop:tmd-css}
（测试用合成陈述：内容从略。）
\uses{prop:tmd-evolution, prop:tmd-small-b}
\end{proposition}
\begin{proofsketch}测试。\end{proofsketch}
\paragraph{地位与缺口}测试：本陈述的地位由\cref{prop:lattice-cs-kernel}决定。

\begin{theorem}[Sivers 函数的符号反转（算符层面） \tagR]\label{thm:tmd-sivers-sign}
（$R_{\mathrm{def}}$：由算符定义直接推出。）
（测试用合成陈述：内容从略。）
非微扰地位：\cref{asm:intro-wightman}、\cref{asm:intro-composite}。
\uses{def:tmd-unsub, def:tmd-collins, asm:intro-wightman, asm:intro-composite}
\end{theorem}

\section{第 6 节（测试）}\label{sec:tmd-evolution}

\begin{corollary}[可观测不对称的符号反转 \tagA]\label{cor:tmd-sivers-observable}
（测试用合成陈述：内容从略。）
非微扰地位：\cref{asm:intro-bridge}。
\uses{thm:tmd-sivers-sign, thm:tmd-factorization, asm:intro-bridge}
\end{corollary}
\begin{proofsketch}测试。\end{proofsketch}
\paragraph{地位与缺口}测试：本陈述的地位由\cref{thm:tmd-factorization}决定。

\begin{proposition}[T-奇碎裂函数的普适性 \tagA]\label{prop:tmd-collins-universality}
（测试用合成陈述：内容从略。）
非微扰地位：\cref{asm:intro-bridge}。
\uses{thm:tmd-factorization, asm:intro-bridge}
\end{proposition}
\begin{proofsketch}测试。\end{proofsketch}
\paragraph{地位与缺口}测试：本陈述的地位由\cref{thm:tmd-factorization}决定。

\begin{definition}[Qiu–Sterman (ETQS) 关联函数]\label{def:tmd-etqs}
（测试用合成陈述：内容从略。）
\uses{def:pdf-quark}
\end{definition}

\begin{proposition}[与扭度 3 共线因子化的匹配 \tagP]\label{prop:tmd-twist3-matching}
（测试用合成陈述：内容从略。）
\uses{def:tmd-etqs, thm:tmd-factorization}
\end{proposition}
\paragraph{地位与缺口}测试：见文献。

\section{第 7 节（测试）}\label{sec:tmd-small-b}

\begin{remark}[$\int d^2k_T\,f\ne f$]\label{rem:tmd-kt-integral}
（测试用合成陈述：内容从略。）
% uses: none
\end{remark}

\begin{example}[单圈全过程]\label{ex:tmd-oneloop}
（测试用合成陈述：内容从略。）
\uses{lem:tmd-soft-oneloop, def:tmd-collins, prop:tmd-evolution-oneloop, prop:tmd-small-b}
\end{example}

\begin{example}[小 $q_T$ 渐近式的重现]\label{ex:tmd-qt-asymptotics}
（测试用合成陈述：内容从略。）
\uses{thm:tmd-factorization, prop:tmd-evolution-oneloop, def:tmd-kin}
\end{example}

\begin{example}[标量双夸克模型中的 Sivers 函数]\label{ex:tmd-sivers-model}
（测试用合成陈述：内容从略。）
\uses{thm:tmd-sivers-sign}
\end{example}
